\documentclass[journal]{IEEEtran}
\usepackage[a5paper, margin=10mm, onecolumn]{geometry}
\usepackage{tfrupee}

\setlength{\headheight}{1cm}
\setlength{\headsep}{0mm}

\usepackage{gvv-book}
\usepackage{gvv}
\usepackage{cite}
\usepackage{amsmath,amssymb,amsfonts,amsthm}
\usepackage{algorithmic}
\usepackage{graphicx}
\usepackage{textcomp}
\usepackage{xcolor}
\usepackage{txfonts}
\usepackage{enumitem}
\usepackage{mathtools}
\usepackage{gensymb}
\usepackage{comment}
\usepackage[breaklinks=true]{hyperref}
\usepackage{tkz-euclide}

\graphicspath{{./figs/}}

\begin{document}
\title{5.3.4}
\author{AI25BTECH11010 - Dhanush Kumar}
\maketitle
\renewcommand{\thefigure}{\theenumi}
\renewcommand{\thetable}{\theenumi}

\textbf{Question}\\

\noindent
Find the product of the matrices 
$\myvec{
1 & 2 & -3 \\
2 & 3 & 2 \\
3 & -3 & -4
}
\myvec{
-6 & 17 & 13 \\
14 & 5 & -8 \\
-15 & 9 & -1
}$
and hence solve the system of linear equations:
$
x + 2y - 3z = -4 \\
2x + 3y + 2z = 2 \\
3x - 3y - 4z = 11
$

\textbf{Solution}\\
\noindent
The given matrices are

\begin{align}
	\vec{A} &= \myvec{1 & 2 & -3 \\ 2 & 3 & 2 \\ -3 & 2 & -4}, &
	\vec{B} &= \myvec{-6 & 14 & -15 \\ 17 & 5 & 9 \\ 13 & -8 & -1}
\end{align}

The product of the matrices is computed as

\begin{align}
	\vec{A}\vec{B} &= \myvec{67 & 0 & 0 \\ 0 & 67 & 0 \\ 0 & 0 & 67} = 67 \vec{I}_3
\end{align}

The given system of linear equations can be written in matrix form as

\begin{align}
	\vec{A}\vec{X} &= \vec{C}, \\
\vec{C} &= \myvec{-4 \\ 2 \\ 1}
\end{align}

Since $\vec{A}\vec{B} = 67 \vec{I}_3$, the inverse of $\vec{A}$ exists and is given by

\begin{align}
	\vec{A}^{-1} = \frac{1}{67}\vec{B}
\end{align}

Multiplying both sides of $A \vec{X} = \vec{C}$ by $\vec{A}^{-1}$, we obtain

\begin{align}
	\vec{X} &= \vec{A}^{-1} \vec{C} = \frac{1}{67} \vec{B} \vec{C}
\end{align}


\begin{align}
	\vec{B} \vec{C} &= \myvec{201 \\ -134 \\ 67}
\end{align}

Therefore, the solution of the system is

\begin{align}
\vec{X} &= \frac{1}{67} \myvec{201 \\ -134 \\ 67} = \myvec{3 \\ -2 \\ 1} \\
\end{align}

\end{document}

